\subsection{片段的 Stokes 公式}

本节中，X 表示类 \(C^{r}\)（\(r \geqslant 2\)）的纯有限维分离流形。A 表示 X 的一个片段，i 表示将 \(\partial A\) 典范嵌入 X 的映射。E 表示 Banach 空间。

11.2.1. 设 \(x \in \partial A\)，且 \(\xi\) 是 \(T_x(\partial A)\) 的一个定向；把 \(\tilde{i}_x(\xi)\) 记作 \(T_x(X)\) 的一个定向，其中含有元素 \(v \wedge u\)，这里 \(v\) 是 A 在 x 处的严格向外向量（11.1.4），而 \(u\) 是 \(\bigwedge^{n-1} T_x(\partial A)\) 中属于定向 \(\xi\) 的非零元素。映射 \(\tilde{i}_x\) 是从 \(\text{Or}(T_x(\partial A))\) 到 \(\text{Or}(T_x(X))\) 的双射。对于映射 \(\tilde{i}_x\)（\(x \in \partial A\)），它定义了一个态射 \(\tilde{i} \colon \partial \widetilde{A} \to \widetilde{X}\)，它是 \(i\) 的一个定向（10.2.5）。若 \(\xi\) 是 X 的一个定向，则 \(\partial A\) 的定向与 \(\xi\) 相关联，并由 \(\tilde{i}\) 得到（10.2.5），称为由 \(\xi\) 定义的定向。

例。--- 当 \(X = R^{n}\) 时，\(\xi\) 是 \(R^{n}\) 的通常定向，且 A 是半径 \(\rho>0\) 的闭球；球面 \(\partial A\) 的定向由 \(\xi\) 定义，就是典范定向（10.2.8, b)）。

11.2.2. 设 \(\omega\) 是 X 上取值于 Banach 空间 E 的 \(p\) 次扭曲微分形式（10.4.1）；其逆像 \(i^{*}(\omega)\) 是 \(\omega\) 通过有向态射 \(i\colon\partial A\to X\) 得到的逆像，记作 \(\omega|\partial A\)，称为 \(\omega\) 在 \(\partial A\) 上诱导的形式（参见 10.4.3）。

11.2.3. 我们作如下两个假设之一：

\begin{enumerate}
\def\labelenumi{(\roman{enumi})}
\tightlist
\item
  \(\omega\) 是 X 上取值于 E 的 \(n - 1\) 次扭曲微分形式；
\item
  X 是有向的，\(\partial A\) 配备与之相应的定向（11.2.1），且 \(\omega\) 是 X 上取值于 E 的 \(n - 1\) 次微分形式。
\end{enumerate}

进一步假设 \(\omega\) 属于类 \(C^{1}\)，且 A 与 \(\omega\) 的支集的交是紧的。\(d\omega\) 是 \(\omega\) 的外微分，并且连续（8.3.5 和 10.3.4）。

A 的示性函数关于 X 上由 \(d\omega\) 定义的向量测度是本质可积的，而 \(\omega|\partial\mathrm{A}\) 是 \(n-1\) 次定义在 \(\partial\mathrm{A}\) 上的微分形式，并且可积（10.4.3 和 10.4.4）。有

\[
\int_ {\mathrm{A}} d \omega = \int_ {\partial \mathrm{A}} \omega \quad (\text{Stokes公式})
\]

11.2.4. 设 \(\alpha\) 是 X 上取值于 Banach 空间 E、具有紧支集、次数为 \(n\) 且类为 \(\mathbf{C}^1\) 的扭曲微分形式。若存在 X 上取值于 E、具有紧支集的扭曲微分形式 \(\omega\)，其次数为 \(n - 1\) 且类为 \(\mathbf{C}^1\)，使得 \(\alpha = d\omega\)，则必须有 \(\int_{\mathbb{X}} \alpha = 0\)。若 X 连通，则此条件充分；此外，若 \(\alpha\) 属于类 \(\mathbf{C}^k\)（\(k\leqslant r-1,\ k\neq\omega\)），则可以选取 \(\omega\) 为类 \(\mathbf{C}^k\)。

11.2.5. 假设 X 是 \(\mathbf{C}\) 的底层实流形，E 是复 Banach 空间，A 是 X 的一个紧片段。赋予 X 由其复结构定义的定向（10.2.7）。设 f 是从 A 到 E 的连续映射，且其在 A 的内部 \(\mathring{A}\) 上的限制是全纯的。以 \(dz\) 表示嵌入 \(z\colon\partial A\to\mathbf{C}\) 的微分。形式 \(f\cdot dz\)，即 f 与 dz 的乘积，是 \(\partial A\) 上取值于 E 的 1 次微分形式，属于类 \(C^{0}\)，并且有：

\[
\int_ {\partial \mathbf {A}} f\, d z = 0 \quad (\text{Cauchy公式})
\]

当 \(f\) 可延拓为定义在 A 的某个开邻域 U 上、仍记作 \(f\) 的全纯映射时，形式 \(f.dz\)（此时 z 表示 U 到 \(\mathbf{C}\) 的典范嵌入）在 U 上属于类 \(C^{\omega}\)，且其外微分为零。

11.2.6（“积分的导数”）。假设 X 是有向的。设 Y 是类 \(C^{r}\) 的流形，\(\alpha\) 是 Y 上取值于 E、具有紧支集、次数为 \(n\)、类为 \(C^{1}\) 的微分形式。设 I 是包含 0 的 R 的开子集，\(g: I \times X \to Y\) 是类 \(C^{r}\) 的态射；对于 \(t \in I\)，记 \(g_{t}\) 为从 X 到 Y 的映射 \(x \mapsto g(t, x)\)。令 \(\psi\) 为从 X 到 T(Y) 的态射，由 \(\psi(x) = T_{(0,x)}(g)(1,0)\) 定义；它是 \(g_{0}\) 的一个提升（8.6.5）。假设 \(\operatorname{pr}_{1}\colon I\times A\to I\) 限制在 \(I\times A\) 与 \(g^{*}(\alpha)\) 的支集的交集上是本原的（TG, I, § 10）。则对于每个 \(t \in I\)，A 与 \(g_{t}^{*}(\alpha)\) 的支集的交是紧的；映射 \(t \mapsto \int_{A} g_{t}^{*}(\alpha)\) 在 I 上属于类 \(C^{1}\)，且其在原点处的导数由公式给出

\[
\frac {d}{d t} \left(\int_ {\mathrm{A}} g _ {t} ^ {*} (\alpha)\right) _ {t = 0} = \int_ {\mathrm{A}} \theta_ {\psi}. \alpha = \int_ {\mathrm{A}} i (\psi) d \alpha + \int_ {\partial \mathrm{A}} i (\psi) \alpha ,
\]

参见第 8.6 号。

